应该学什么，应该教什么?
与人类的记忆一样，人工神经网络的记忆在可靠性方面也没有保障，无论它的编码是在神经网络的连接之中还是在网络环路传播的数据之中。我们身边的计算机至少在两项计算任务上远远超越了我们：计算速度以及信息的可靠储存。我敢打赌，明天的人工智能会知道怎么利用这些实实在在超越人类的能力，它们可能将神经网络的某些方面与其他算法结合起来，这些算法更可靠，也对更特殊的任务进行了优化。

新技术可能已经改变了你的大脑皮层。我们沉迷于智能手机、谷歌和维基百科，这似乎影响了我们对记忆的管理。这不一定是坏事。

我认为，知识和技能的过剩损害了对概念和模型的理解，尤其是对理解这个世界来说最有用、最可信的概念和模型。这种教育往往忽略了模型可信的原因及其适用范围。我的意见是，教授的内容应该大量减少，教师应该只教授那些违反直觉并且有教育意义的重要内容。比如说，我认为应该教授认知偏差、演化理论的关键过程、理论计算机科学和道德功利主义，同时可以削减三角学和量子力学等内容。

此外，贝叶斯公式似乎提示我们应该通过例子来学习，而不是直接记住理论——我们会在接下来的章节中看到，我们的大脑似乎偏向贝叶斯主义，从具体事例出发很快就能推广到一般情况。因此，似乎只有在获得了足够的数据，并使这些数据可以“轻松访问”，让我们能够轻易估计贝叶斯公式中的思想实验项之后，理论的重要性才会突然显现。这样的话，似乎在考虑大量例子，使理论的**效用**凸显眼前之后，再去学习这些理论才更合适。比如说，我们应该先以游戏、谜题和逻辑悖论等能吸引学生的形式来引入数学，之后再向他们解释这些内容都是更普遍的理论的应用事例——这正是我在这本书中努力尝试的做法。

但依我愚见，要教授的最重要的内容还是认识论，还有对于认识论的应用来说不可或缺的统计学。当然，作为极端贝叶斯主义者，我尤其认为贝叶斯公式及其大量违反直觉的推论应该成为教育的支柱之一。
我认为是时候放弃积累公认正确的教条知识这种做法了，应当转而教授知识是什么、如何获得知识、如何分辨可信的理论和不值得赋予置信度的理论。不幸的是，在我们这个时代，即使是大科学家也非常欠缺对认识论的理解，很多人甚至不知道贝叶斯主义的存在。